\documentclass[11pt,a4paper]{article}

% ============== PACKAGES ==============
\usepackage[utf8]{inputenc}
\usepackage[T1]{fontenc}
\usepackage[french]{babel}
\usepackage[margin=2.3cm]{geometry}
\usepackage{graphicx}
\usepackage[table]{xcolor}
\usepackage{titlesec}
\usepackage{fancyhdr}
\usepackage{array}
\usepackage{tabularx}
\usepackage{booktabs}
\usepackage{enumitem}
\usepackage{amssymb}
\usepackage{listings}
\usepackage{hyperref}
\usepackage{tikz}
\usetikzlibrary{shapes.geometric, arrows.meta, positioning}

% ============== COLORS ==============
\definecolor{primaryblue}{RGB}{31, 78, 140}
\definecolor{lightblue}{RGB}{220, 232, 248}
\definecolor{codebg}{RGB}{245, 245, 245}
\definecolor{codegreen}{RGB}{0, 128, 0}
\definecolor{codeorange}{RGB}{215, 95, 0}
\definecolor{darkgray}{RGB}{80, 80, 80}

% ============== HYPERREF ==============
\hypersetup{
    colorlinks=true,
    linkcolor=primaryblue,
    urlcolor=primaryblue,
    citecolor=primaryblue
}

% ============== TITLE FORMATTING ==============
\titleformat{\section}
    {\color{primaryblue}\Large\bfseries}
    {\thesection.}{0.5em}{}
    [\vspace{-0.3em}\color{primaryblue}\rule{\linewidth}{0.6pt}]

\titleformat{\subsection}
    {\color{primaryblue}\large\bfseries}
    {\thesubsection}{0.5em}{}

% ============== HEADERS / FOOTERS ==============
\pagestyle{fancy}
\fancyhf{}
\fancyhead[L]{\small\color{darkgray}Compte Rendu --- Microservices Résilients}
\fancyhead[R]{\small\color{darkgray}INGA2 G1}
\fancyfoot[L]{\small\color{darkgray}\textit{Votre Nom} --- ING252 Architecture Logicielle}
\fancyfoot[R]{\small\color{darkgray}Page \thepage}
\renewcommand{\headrulewidth}{0.4pt}
\renewcommand{\footrulewidth}{0.4pt}

% ============== CODE LISTINGS ==============
\lstdefinestyle{javastyle}{
    backgroundcolor=\color{codebg},
    basicstyle=\ttfamily\footnotesize,
    keywordstyle=\color{primaryblue}\bfseries,
    commentstyle=\color{codegreen}\itshape,
    stringstyle=\color{codeorange},
    numbers=left,
    numberstyle=\tiny\color{darkgray},
    numbersep=8pt,
    frame=single,
    rulecolor=\color{darkgray},
    breaklines=true,
    showstringspaces=false,
    tabsize=2,
    captionpos=b
}

\lstdefinestyle{yamlstyle}{
    backgroundcolor=\color{codebg},
    basicstyle=\ttfamily\footnotesize,
    keywordstyle=\color{primaryblue}\bfseries,
    commentstyle=\color{codegreen}\itshape,
    numbers=left,
    numberstyle=\tiny\color{darkgray},
    numbersep=8pt,
    frame=single,
    breaklines=true,
    showstringspaces=false,
    tabsize=2
}

\lstset{style=javastyle}

% ============== TABLE STYLE ==============
\renewcommand{\arraystretch}{1.3}
% Pour utiliser \rowcolor sur la ligne d'en-tete d'une tabularx, il faut une astuce :
\newcommand{\hdrcell}[1]{\textbf{\textcolor{white}{#1}}}

% ============== DOCUMENT ==============
\begin{document}

% ===================== TITLE PAGE / HEADER =====================
\begin{center}
    \vspace*{1cm}
    {\color{primaryblue}\Huge\bfseries COMPTE RENDU}\\[0.5em]
    {\color{primaryblue}\Large\bfseries Architecture Microservices Résilients}\\[0.5em]
    {\large Circuit Breaker \& Bulkhead avec Spring Boot}\\[1em]
    \color{primaryblue}\rule{0.7\linewidth}{1.2pt}\\
\end{center}

\vspace{0.5cm}

\begin{center}
\begin{tabular}{rl}
    \textbf{Auteur} & : Votre Nom \\
    \textbf{Matière} & : ING252 --- Architecture Logicielle \\
    \textbf{Classe} & : INGA2 --- Groupe 1 \\
    \textbf{Date} & : Mai 2026 \\
    \textbf{Repository} & : \texttt{github.com/votre-username/microservices-resilience} \\
\end{tabular}
\end{center}

\vspace{1cm}

% ===================== 1. INTRODUCTION =====================
\section{Introduction}

Ce TP a pour objectif d'implémenter une architecture microservices simple en Spring Boot, intégrant deux patterns de résilience essentiels : le \textbf{Circuit Breaker} et le \textbf{Bulkhead}. Ces patterns permettent de garantir la robustesse du système face aux défaillances et à la surcharge des services distants.

\subsection*{Patterns implémentés}
\begin{itemize}[leftmargin=*, label=\textcolor{primaryblue}{$\blacktriangleright$}]
    \item \textbf{Circuit Breaker} (Resilience4j) --- protection contre les défaillances en cascade
    \item \textbf{Bulkhead Pattern} (Resilience4j) --- isolation des ressources et limitation de la concurrence
    \item \textbf{Fallback} --- réponses dégradées en cas d'échec
\end{itemize}

\subsection*{Technologies utilisées}
\begin{itemize}[leftmargin=*, label=\textcolor{primaryblue}{$\blacktriangleright$}]
    \item Java 17, Spring Boot 3.2
    \item Resilience4j 2.2 (\texttt{resilience4j-spring-boot3})
    \item Maven (projet multi-modules)
    \item Spring Boot Actuator (monitoring des patterns)
\end{itemize}

% ===================== 2. ARCHITECTURE =====================
\section{Architecture Globale}

L'architecture est composée de trois microservices Spring Boot indépendants. Le \textit{Order Service} joue le rôle de consommateur et applique les patterns de résilience sur ses appels vers les deux autres services.

\subsection*{Schéma de l'architecture}

\begin{center}
\begin{tikzpicture}[
    node distance=1.5cm and 2cm,
    service/.style={rectangle, rounded corners, draw=primaryblue, fill=lightblue, thick, minimum width=3.2cm, minimum height=1.1cm, align=center},
    arr/.style={-Stealth, thick, primaryblue},
    patternlbl/.style={font=\scriptsize\bfseries, color=codeorange, fill=yellow!20, rounded corners, inner sep=2pt}
]
    \node[service] (client) {Client\\\small(test.ps1 / curl)};
    \node[service, right=3.5cm of client] (order) {Order Service\\Port 8080};

    \node[service, above right=0.8cm and 3.5cm of order] (inventory) {Inventory Service\\Port 8081};
    \node[service, below right=0.8cm and 3.5cm of order] (payment) {Payment Service\\Port 8082};

    \draw[arr] (client) -- (order);
    \draw[arr] (order.east) -- (inventory.west)
        node[patternlbl, midway, above, sloped] {Circuit Breaker};
    \draw[arr] (order.east) -- (payment.west)
        node[patternlbl, midway, below, sloped] {Bulkhead};
\end{tikzpicture}
\end{center}

\subsection*{Rôle de chaque composant}

\begin{tabularx}{\linewidth}{|l|c|X|}
    \hline
    \rowcolor{primaryblue} \hdrcell{Service} & \hdrcell{Port} & \hdrcell{Rôle} \\
    \hline
    Order Service & 8080 & Service consommateur. Applique les patterns Circuit Breaker et Bulkhead sur ses appels distants. \\
    \hline
    Inventory Service & 8081 & Service métier simulant des défaillances aléatoires (60\%) pour tester le Circuit Breaker. \\
    \hline
    Payment Service & 8082 & Service de paiement avec un traitement lent (3s) pour tester le Bulkhead. \\
    \hline
\end{tabularx}

\subsection*{Structure du projet}

\begin{lstlisting}[basicstyle=\ttfamily\footnotesize, frame=single, backgroundcolor=\color{codebg}]
microservices-resilience/
|-- pom.xml                      (parent)
|-- order-service/               (Circuit Breaker + Bulkhead)
|   |-- pom.xml
|   `-- src/main/
|       |-- java/com/tp/order/
|       |   |-- OrderServiceApplication.java
|       |   |-- OrderController.java
|       |   |-- InventoryClient.java
|       |   `-- PaymentClient.java
|       `-- resources/application.yml
|-- inventory-service/           (echecs aleatoires)
|   `-- src/main/java/com/tp/inventory/
|       |-- InventoryServiceApplication.java
|       `-- InventoryController.java
|-- payment-service/             (traitement lent)
|   `-- src/main/java/com/tp/payment/
|       |-- PaymentServiceApplication.java
|       `-- PaymentController.java
`-- test.ps1                     (script de test)
\end{lstlisting}

% ===================== 3. CIRCUIT BREAKER =====================
\section{Circuit Breaker (Inventory Service)}

\subsection*{Principe}

Le Circuit Breaker protège le système en détectant les défaillances et en coupant les appels vers un service défaillant. Cela évite la propagation des erreurs et permet au service défaillant de récupérer.

\subsection*{États du Circuit Breaker}

\begin{tabularx}{\linewidth}{|l|X|}
    \hline
    \rowcolor{primaryblue} \hdrcell{État} & \hdrcell{Description} \\
    \hline
    \texttt{CLOSED} & Fonctionnement normal --- les appels passent vers le service. \\
    \hline
    \texttt{OPEN} & Circuit ouvert --- tous les appels sont rejetés immédiatement (Fast-Fail). \\
    \hline
    \texttt{HALF\_OPEN} & Test de récupération avec un nombre limité d'appels. \\
    \hline
\end{tabularx}

\subsection*{Configuration (\texttt{application.yml})}

\begin{lstlisting}[style=yamlstyle, language=, caption={Configuration Circuit Breaker}]
resilience4j:
  circuitbreaker:
    instances:
      inventoryService:
        sliding-window-size: 10
        minimum-number-of-calls: 5
        failure-rate-threshold: 50
        wait-duration-in-open-state: 10s
        permitted-number-of-calls-in-half-open-state: 3
\end{lstlisting}

\subsection*{Implémentation Java}

L'annotation \texttt{@CircuitBreaker} de Resilience4j enveloppe la méthode et déclenche le fallback en cas d'échec ou de circuit ouvert.

\begin{lstlisting}[language=Java, caption={InventoryClient.java}]
@Service
public class InventoryClient {

    private final RestTemplate restTemplate;

    @Value("${services.inventory.url}")
    private String inventoryUrl;

    public InventoryClient(RestTemplate restTemplate) {
        this.restTemplate = restTemplate;
    }

    @CircuitBreaker(name = "inventoryService",
                    fallbackMethod = "checkStockFallback")
    public String checkStock(String productId) {
        String url = inventoryUrl + "/inventory/" + productId;
        return restTemplate.getForObject(url, String.class);
    }

    private String checkStockFallback(String productId, Throwable t) {
        return "FALLBACK: Inventory service unavailable for product "
               + productId;
    }
}
\end{lstlisting}

\subsection*{Comportement attendu}
Le service \textit{Inventory} simule des échecs aléatoires (60\% de chance). Après 5 échecs dans la fenêtre de 10 appels, le circuit s'ouvre et les appels suivants sont rejetés immédiatement sans atteindre le service.

% ===================== 4. BULKHEAD =====================
\section{Bulkhead Pattern (Payment Service)}

\subsection*{Principe}

Le Bulkhead isole les ressources en limitant le nombre d'appels simultanés vers un service. Il empêche un service lent de monopoliser tous les threads disponibles, ce qui pourrait dégrader l'ensemble de l'application.

\subsection*{Configuration}

\begin{itemize}[leftmargin=*, label=\textcolor{primaryblue}{$\blacktriangleright$}]
    \item \textbf{Type} : Semaphore (isolation au niveau des threads)
    \item \textbf{Max appels concurrents} : 2
    \item \textbf{Comportement si plein} : Rejet immédiat (\texttt{BulkheadFullException}) puis fallback
\end{itemize}

\begin{lstlisting}[style=yamlstyle, language=, caption={Configuration Bulkhead}]
resilience4j:
  bulkhead:
    instances:
      paymentService:
        max-concurrent-calls: 2
        max-wait-duration: 0
\end{lstlisting}

\subsection*{Implémentation Java}

\begin{lstlisting}[language=Java, caption={PaymentClient.java}]
@Service
public class PaymentClient {

    private final RestTemplate restTemplate;

    @Value("${services.payment.url}")
    private String paymentUrl;

    public PaymentClient(RestTemplate restTemplate) {
        this.restTemplate = restTemplate;
    }

    @Bulkhead(name = "paymentService",
              fallbackMethod = "processPaymentFallback")
    public String processPayment(String orderId) {
        String url = paymentUrl + "/payment/" + orderId;
        return restTemplate.getForObject(url, String.class);
    }

    private String processPaymentFallback(String orderId, Throwable t) {
        return "REJECTED: Bulkhead full - payment service overloaded "
               + "for order " + orderId;
    }
}
\end{lstlisting}

\subsection*{Comportement attendu}
Lors de l'envoi de 5 requêtes simultanées vers \texttt{/orders/pay}, seules 2 sont traitées en parallèle. Les 3 autres sont rejetées immédiatement avec le message \og Bulkhead full \fg.

% ===================== 5. SERVICES SIMULES =====================
\section{Services Simulés}

\subsection*{Inventory Service (échecs aléatoires)}

\begin{lstlisting}[language=Java, caption={InventoryController.java}]
@RestController
@RequestMapping("/inventory")
public class InventoryController {

    private final Random random = new Random();

    @GetMapping("/{productId}")
    public ResponseEntity<String> getStock(@PathVariable String productId) {
        // 60% de chance d'echec pour declencher le Circuit Breaker
        if (random.nextInt(100) < 60) {
            return ResponseEntity
                .status(HttpStatus.INTERNAL_SERVER_ERROR)
                .body("Simulated failure for product " + productId);
        }
        int stock = 10 + random.nextInt(50);
        return ResponseEntity.ok("Product " + productId + " - Stock: " + stock);
    }
}
\end{lstlisting}

\subsection*{Payment Service (traitement lent)}

\begin{lstlisting}[language=Java, caption={PaymentController.java}]
@RestController
@RequestMapping("/payment")
public class PaymentController {

    @GetMapping("/{orderId}")
    public String processPayment(@PathVariable String orderId)
            throws InterruptedException {
        // Traitement lent simule : 3 secondes
        Thread.sleep(3000);
        return "Payment processed for order " + orderId;
    }
}
\end{lstlisting}

% ===================== 6. TESTS =====================
\section{Tests et Résultats}

Le script PowerShell \texttt{test.ps1} (disponible dans le repository) automatise l'ensemble des tests.

\subsection{Test Circuit Breaker}

\textbf{Scénario} : 12 appels consécutifs au Service Inventory avec simulation d'échecs aléatoires.

\begin{tabularx}{\linewidth}{|X|c|}
    \hline
    \rowcolor{primaryblue} \hdrcell{Métrique} & \hdrcell{Valeur} \\
    \hline
    Succès & 2 \\
    \hline
    Échecs (erreurs simulées) & 5 \\
    \hline
    Fallback (circuit ouvert) & 5 \\
    \hline
\end{tabularx}

\textbf{Analyse} : Après 5 échecs dans la fenêtre de 10 appels (50\%), le circuit passe à l'état \texttt{OPEN}. Les appels suivants sont rejetés immédiatement par le Circuit Breaker sans appeler le service distant (Fast-Fail).~\checkmark

\subsection{Test Bulkhead}

\textbf{Scénario} : 5 requêtes de paiement envoyées simultanément.

\begin{tabularx}{\linewidth}{|X|c|}
    \hline
    \rowcolor{primaryblue} \hdrcell{Métrique} & \hdrcell{Valeur} \\
    \hline
    Acceptés & 2 \\
    \hline
    Rejetés (Bulkhead plein) & 3 \\
    \hline
\end{tabularx}

\textbf{Analyse} : Avec une limite de 2 appels concurrents, seuls 2 paiements sont traités. Les 3 autres sont immédiatement rejetés et passent par le fallback.~\checkmark

\subsection{Résumé des tests}

\begin{tabularx}{\linewidth}{|X|X|c|}
    \hline
    \rowcolor{primaryblue} \hdrcell{Test} & \hdrcell{Résultat} & \hdrcell{Statut} \\
    \hline
    Circuit Breaker & Circuit s'ouvre après 5 échecs & \checkmark \\
    \hline
    Fallback Circuit Breaker & Réponse dégradée retournée & \checkmark \\
    \hline
    Bulkhead & Max 2 appels, rejet des suivants & \checkmark \\
    \hline
    Fallback Bulkhead & Message \og Bulkhead full \fg{} retourné & \checkmark \\
    \hline
\end{tabularx}

\subsection*{Endpoints de monitoring (Actuator)}

Spring Boot Actuator expose l'état des patterns en temps réel :

\begin{itemize}[leftmargin=*, label=\textcolor{primaryblue}{$\blacktriangleright$}]
    \item \texttt{GET /actuator/health} --- état de santé global
    \item \texttt{GET /actuator/circuitbreakers} --- état du Circuit Breaker (CLOSED / OPEN / HALF\_OPEN)
    \item \texttt{GET /actuator/bulkheads} --- nombre d'appels disponibles dans le Bulkhead
\end{itemize}

% ===================== 7. CONCLUSION =====================
\section{Conclusion}

Ce TP a permis de mettre en œuvre concrètement deux patterns de résilience essentiels dans une architecture microservices Spring Boot. Les deux scénarios de défaillance --- échecs aléatoires et surcharge --- ont été reproduits et validés expérimentalement.

\subsection*{Points retenus}
\begin{itemize}[leftmargin=*, label=\textcolor{primaryblue}{$\blacktriangleright$}]
    \item Le \textbf{Circuit Breaker} protège efficacement contre les défaillances en cascade en coupant rapidement les appels vers un service en échec.
    \item Le \textbf{Bulkhead} protège les ressources locales en limitant la concurrence sur les services distants potentiellement lents.
    \item Le mécanisme de \textbf{fallback} permet de fournir des réponses dégradées plutôt que des erreurs brutes, améliorant l'expérience utilisateur.
    \item L'utilisation de \textbf{Resilience4j} via les annotations rend l'intégration des patterns simple et déclarative.
\end{itemize}

\subsection*{Tableau de validation}

\begin{tabularx}{\linewidth}{|l|l|X|}
    \hline
    \rowcolor{primaryblue} \hdrcell{Pattern} & \hdrcell{Technologie} & \hdrcell{Validation} \\
    \hline
    Circuit Breaker & Resilience4j & Circuit OPEN après 5 échecs \checkmark \\
    \hline
    Bulkhead & Resilience4j & Rejet après 2 appels simultanés \checkmark \\
    \hline
    Fallback & Resilience4j & Réponse dégradée retournée \checkmark \\
    \hline
    Monitoring & Spring Actuator & Métriques exposées \checkmark \\
    \hline
\end{tabularx}

\vspace{0.5cm}

Ce travail démontre la capacité à concevoir un système distribué tolérant aux pannes, répondant aux exigences de disponibilité et de robustesse des architectures microservices modernes.

\end{document}
